\chapter{Introduction}
\label{ch:intro}

%-------------------------------------------------------

Trading card collectors keep track of thousands of cards, and the value of a card depends on its exact version: the set it was printed in, its number inside the set, its rarity, its printed language and its physical condition. Typing this information by hand is slow and error-prone. The goal of this project is to identify the exact card from an ordinary photo taken with a phone, and to build a small application around that capability.

\section{Problem statement}

Given a photo that contains one (or several) Pok\'emon trading cards, the system must return, for each card, the identifier of the matching card in a reference database of official card scans (for example \texttt{sv9-116}: set \emph{Journey Together}, number 116). Formally, this is \emph{instance-level image retrieval}: the query is the photographed card and the database contains exactly one reference image per card id. The output is a full ranking of the database, so that both top-1 and top-$k$ accuracy can be measured.

Three properties make the problem harder than it first appears:
\begin{enumerate}
  \item \textbf{Near-duplicates and reprints.} A new set can reprint a card of an older set with identical artwork, and many Pok\'emon appear in several versions (regular, illustration rare, special illustration rare). In the reference database used here, \res{db-art-pairs} pairs of cards share identical artwork and \res{db-dup-names} card names occur more than once (\res{db-dup-name-cards} cards in total; \emph{\res{db-max-name}} alone has \res{db-max-name-n} versions). A method that recognises only ``which Pok\'emon'' is therefore not enough.
  \item \textbf{Photo conditions.} Phone photos show the card in perspective, rotated, possibly upside down, with uneven lighting, glare on glossy or holographic surfaces, motion blur and sensor noise.
  \item \textbf{Text and language.} The same card is printed in several languages. The artwork is identical, but the text (name, attacks, rules) differs, so a method that relies on the text area can fail on non-English cards.
\end{enumerate}

\section{Motivation}

Instance retrieval of flat, textured objects is a classical computer-vision problem, and trading cards are a good test bed to compare the main families of methods on the same data: hand-crafted local features with geometric verification \citep{lowe2004sift, rublee2011orb, fischler1981ransac}, global embeddings from a convolutional network trained on ImageNet \citep{he2016resnet}, and embeddings from a vision transformer trained with language supervision \citep{dosovitskiy2021vit, radford2021clip}. The comparison is interesting because the methods have opposite strengths: local features are precise but slow when matched exhaustively, while global embeddings are fast but may confuse visually similar cards. A hybrid that uses the embedding to build a short candidate list and local features to verify it is a natural compromise, and the project measures how much of each strength it keeps.

\section{Contributions}

The contributions of this project are:
\begin{itemize}
  \item a reference database of \res{db-cards} cards from \res{db-sets} recent sets, with reprints grouped by artwork using perceptual hashing, and two evaluation metrics (strict and lenient) that handle reprints explicitly (Chapter~\ref{ch:data});
  \item a classical card detection and rectification stage, and a common retrieval interface implemented for seven methods and their variants (Chapters~\ref{ch:approach} and~\ref{ch:methods});
  \item a synthetic query generator with five distortion severities and exact ground truth, and an evaluation that logs every query with 95\% bootstrap confidence intervals;
  \item a systematic comparison of all methods across the five severities, a study of the hybrid's shortlist size, a rectification ablation, multi-card scenes, a cross-language test and an automatic failure analysis (Chapters~\ref{ch:results} and~\ref{ch:discussion});
  \item two extensions with their own validation: printed-language recognition with OCR and a visual condition estimate;
  \item a working local web application (CardDex) that integrates all of the above, recognises English and Japanese printings, and was used to identify and catalogue real cards; the photos saved in it serve as a small real-photo test of all methods (Section~\ref{sec:res-app}).
\end{itemize}
